\section{Discussion}
\label{sec:scope-transfer}

The local matching theorem identifies how a perturbation invisible to the
stationary projection can alter a scalar fold condition.  The first delay
moment produces transverse forcing, the stable transverse equation gives
the leading coefficient \(A_N^{-1}\mathsf S_N\eta\), and the heterogeneous quadratic
terms return this response to the collective solvability condition through
\(\mathfrak r_N\).  The resulting functional
\(\Lambda_N=\mathfrak r_NA_N^{-1}\mathsf S_N\) is realized by derivatives
of a matching function on compatible RFDE histories.  Equality of projected
vector fields at the same history is therefore insufficient to determine
the matching sensitivity.

The full-range statement concerns the unrestricted matrix tangent space
\(\mathcal T_N\).  Under a prescribed pattern, the available forcing is
exactly the image of \(R\mapsto R\mathbf1\) in
\cref{thm:rw-structured-delay-range}.  If every delay layer preserves its
row sums, then \(\mathsf S_N=0\) and the coefficient computed here
vanishes.  This identifies a degeneracy of the first-moment contribution;
it gives no conclusion about higher-order sensitivities.  Likewise, the
vanishing of this map when only one delay remains refers to the stated
constraints, including preservation of the layer sum.

Uniformity in network size uses more than invertibility of each finite
matrix \(A_N\).  Dobrushin contraction supplies both a common inverse
bound for the coefficient and common exponential decay for the invariant
history construction.  The mixing family in
\cref{prop:rw-growing-families} has a nonzero response for all \(N\),
whereas the cyclic family loses the uniform inverse bound.  The present
theorem consequently does not cover arbitrary sequences of increasingly
large local networks.  A treatment of such families would have to track
the interaction between the deteriorating transverse decay and
\(\delta\to0\).

For another RFDE, the algebraic elimination survives whenever a bounded
transverse inverse and a collective operator with one-dimensional cokernel
are available and the parameter direction is transverse to the collective
range.  Its realization as a local matching derivative requires a
compatible-history invariant manifold and uniform differentiated matching
estimates.  The affine fold orbit, Gaussian cokernel, and explicit pairing
with the quadratic coefficients are features of the present polynomial
model; they must be replaced by the corresponding calculations in a new
equation.  A global connection would further require the invariant-manifold
defining function and the comparison bounds in
\cref{cor:rw-conditional-connection}.  These distinctions locate the
additional mathematical work needed to pass from the local coefficient to
a different model or to a global parameter locus.
